\section{Đường thứ hai, và cái giá của nó không phải giá tiền}
\label{sec:ch16-duong-thu-hai}

\index{đánh giá!đường thứ hai trả lời câu nào|(}%
Câu mà đường thứ nhất trả lời thì mục trước đã nêu, và nó đúng là câu của
định nghĩa~\ref{def:ch01-do-trung-thuc}. Đường thứ hai trả lời một câu khác, và
mục~\ref{sec:ch12-hai-nhanh-danh-gia} đã dựng sẵn phép phân biệt cần dùng, nên
mục này chỉ phải phát biểu lại nó ở mức một lối thoát chứ không ở mức một
chương.

\index{đánh giá!người đọc không nhìn thấy mô hình}%
Phép phân biệt ấy đứng trên một sự thật đơn giản về vị trí: người đọc không nhìn
thấy bên trong mô hình. Cho nên dù hỏi bao nhiêu người, câu trả lời của họ nói
về quan hệ giữa lời giải thích và người đọc, tức
\tn{tính hợp lý}{plausibility} theo dòng thứ ba của
bảng~\ref{tab:ch01-bon-quan-he}, chứ không nói về quan hệ giữa lời giải thích và
mô hình. Một chỉ số trung thực hoàn hảo vẫn có thể không tương quan với cảm nhận
của người, và điều đó không phải bằng chứng chống lại nó.

Trong phạm vi ấy thì đường thứ hai vẫn xác nhận được đúng một loại lời đề xuất,
và loại ấy có thật. Khi một chỉ số tự động được đề xuất như
\tn{đại lượng thay thế}{proxy} cho chất lượng của lời giải thích
\emph{đối với người dùng}, thì phép thử tối thiểu là cho thấy điểm số của nó đi
cùng đánh giá của người trên cùng những lời giải thích ấy. Đó là món nợ mà
mục~\ref{sec:ch12-hai-nhanh-danh-gia} đi đòi, và nó là món nợ có thật, vì chính
Doshi-Velez và Kim đặt điều kiện ấy lên cách đánh giá theo chức
năng~\autocite{p14rigorous}.

\index{kiểm đếm tài liệu XAI!đọc như một lối thoát}%
Phép thử ấy đã được chạy bao nhiêu lần thì
mục~\ref{sec:ch12-con-so-va-cai-no-dem} đã đếm: trong 18\,254 bài mang từ khóa
về \tn{khả năng giải thích}{explainability}, 253 bài có từ khóa về thử nghiệm
với người, 237 bài đúng
chủ đề, và 128 bài thật sự có kiểm chứng với người, tức
0.70\%~\autocite{p24humangap}. Đọc con số ấy như một phát biểu về lối thoát chứ
không như một phát biểu về một chương thì nó nói rằng đường thứ hai tồn tại, đi
được, và gần như không ai đi.

\index{đánh giá!trần của đường thứ hai, theo lời bài báo}%
Chỗ mà mục này thêm vào so với chương~\ref{ch:con-nguoi-vang-mat} là trần của
đường ấy nói bằng chính chữ của bài báo, vì bài báo phát biểu nó hai lần và cả
hai lần đều chọn danh từ rất cẩn thận. Ở phần kết, trang 6: không có bằng chứng
thực nghiệm thì không thể đánh giá được liệu các phương pháp giải thích có thật
sự đáp ứng \emph{nhu cầu của người dùng cuối} hay không. Và ở mục 5.2, cùng
trang: kiểm thử thực nghiệm với người nhất thiết là chuẩn vàng cho
\emph{những lời tuyên bố về khả năng giải thích}~\autocite{p24humangap}.

\index{độ trung thực!vắng tên trong bài của đường thứ hai}%
Hai danh từ ấy, nhu cầu của người dùng cuối và lời tuyên bố về khả năng giải
thích, không danh từ nào là đối tượng của định
nghĩa~\ref{def:ch01-do-trung-thuc}. Điều đó khớp với phép rà toàn văn mà
mục~\ref{sec:ch12-con-so-va-cai-no-dem} đã ghi và mục này kiểm lại: trong bảy
trang của bài, các chữ \enquote{faithful}, \enquote{faithfulness},
\enquote{ground truth} và \enquote{fidelity} không xuất hiện lần
nào~\autocite{p24humangap}. Bài không tránh né chủ đề độ trung thực; nó không ở
trong chủ đề ấy. Đọc theo bảng~\ref{tab:ch01-bon-quan-he} thì nó cấp bằng chứng
cho dòng tính hợp lý và không cấp cho dòng độ trung thực, và cột cuối của chính
bảng ấy ghi rằng từ dòng thứ ba không suy ra được lời giải thích đúng với mô
hình.

\index{đánh giá!chi phí không phải chỗ tắc}%
Còn cái giá thì đây là chỗ chương phải sửa lại câu hỏi của chính mình. Cách đọc
thông thường về việc chỉ 128 bài trên 18\,254 chạy kiểm chứng với người là cách
đọc theo chi phí: thí nghiệm có người thì đắt, chậm, phải tuyển người, phải qua
thủ tục, nên người ta bỏ qua. Cách đọc ấy không lấy được từ bài báo. Rà toàn văn
bảy trang thì các chữ \enquote{expensive}, \enquote{power}, \enquote{IRB},
\enquote{consent}, \enquote{budget} và \enquote{burden} không xuất hiện lần nào,
còn chữ \enquote{cost} xuất hiện đúng một lần và nằm trong tên một tác giả ở
thư mục tham khảo~\autocite{p24humangap}. Bài không phân tích chi phí, không ước
lượng chi phí, và không nêu chi phí như lý do.

Câu duy nhất trong bài đọc gần giống một cái giá nằm ở trang 5, và nó nói về
việc tập hợp chuyên gia: tất nhiên có thể khó tập hợp được các chuyên gia lĩnh
vực để thử phương pháp và công cụ, nhất là những chuyên gia mà ta không cùng làm
công cụ với họ, và có lẽ vì thế mà một số bài chỉ có hai hoặc ba chuyên gia
trong phần đánh giá. Câu ngay sau đó rút lại lời nhân nhượng vừa nêu: dù vậy, khó mà
tin những phát biểu có tính khái quát về khả năng giải thích khi nhóm người tham
gia rất nhỏ hoặc không giống người dùng đích~\autocite{p24humangap}.

\index{đánh giá!bài báo nêu phản đối rồi bác}%
Ở trang 6 thì bài nêu hai lời phản đối thông thường với đánh giá thực nghiệm và
bác cả hai, và đoạn ấy quyết định cách mục này phải viết. Hai lời phản đối mà
bài nêu không phải lời phản đối về chi phí: chúng là con người thì thiên lệch,
và khó tập hợp được một mẫu người đại diện. Câu trả lời của bài thì đi qua đúng
chỗ chi phí. Bài viết rằng cũng có thể nói y như vậy về mọi benchmark học máy
đang dùng, từ ImageNet tới các bộ đánh giá bình luận phim và hỏi đáp y khoa: tất
cả đều mang thiên lệch của con người từ khâu sinh ra, chọn lọc và bối cảnh văn
hóa, có thể khớp hoặc không khớp với ứng dụng đích, và \emph{đã tốn rất nhiều
thời gian để thu thập; dù vậy người ta vẫn dùng chúng}. Rồi bài kết luận rằng
các phương pháp để thử XAI thì rõ ràng đã có, và tuy có những khó khăn khi
chuyển từ chỗ chỉ thử AI trên các bộ dữ liệu tĩnh sang chỗ thử với người,
\emph{những khó khăn ấy không phải là không vượt qua nổi}~\autocite{p24humangap}.
Nghĩa là công sức bỏ ra được bài nêu đúng một lần, ở vế bị so sánh, và nêu để
nói rằng ngành vẫn trả nó cho những thứ khác.

\index{đánh giá!công trình mà bài dẫn để nói phương pháp đã có}%
Công trình mà bài dẫn ở chỗ nói các phương pháp đã có là bài của Doshi-Velez và
Kim~\autocite{p14rigorous}, tức đúng bài mà
mục~\ref{sec:ch01-danh-gia-giai-thich} lấy ba tầng đánh giá và
mục~\ref{sec:ch12-hai-nhanh-danh-gia} lấy quan hệ vay mượn giữa chúng. Điều bài
nói ở đó chỉ là phương pháp đã có sẵn. Bước từ đó sang câu \enquote{vậy chỗ tắc
không phải phương pháp mà là việc không ai bị đòi phải dùng nó} là bước của
sách, không của bài; cái bài đưa ra để đỡ cho bước ấy là một câu nhắn gửi người
bình duyệt, rằng có trách nhiệm đẩy lại những nghiên cứu tuyên bố giải thích
được mà đưa rất ít hoặc không đưa bằng chứng
nào~\autocite{p24humangap}. Một lời nhắn gửi người bình duyệt là một phát biểu về
chỗ nào trong quy trình chưa đòi hỏi điều gì, nên nó đỡ được bước ấy mà không
chứng minh nó.

\index{đánh giá!cái giá thật là chất lượng bằng chứng}%
Bài lại có ghi một cái giá, và nó ở chỗ khác hẳn: không phải giá phải trả để
chạy phép thử mà là chất lượng của bằng chứng thu được sau khi chạy. Trang 5:
trong 128 bài có kiểm chứng thì 13 bài không báo số người tham gia và 3 bài chỉ
báo con số ước chừng, và ở nhiều bài thì các tác giả không nêu chuyên môn, nơi
làm việc, học vấn hay quan hệ của người tham gia với chính
họ~\autocite{p24humangap}. Cộng với quan sát vừa dẫn về những phần đánh giá chỉ
có hai hoặc ba chuyên gia, điều bài mô tả không phải một lối thoát chưa ai đủ
tiền đi mà một lối thoát đã có người đi và đi ẩu.

\index{đánh giá!người tự tin không phải người đúng}%
Chỗ sắc nhất trong phần ấy là chỗ bài cho thấy vì sao chất lượng lại quan trọng
đến thế, và nó là một câu về khoảng cách giữa tự tin và đúng. Bài viết rằng một
phép xác nhận rằng tín hiệu đã được nhận và được hiểu vẫn là cần thiết, vì công
trình trước đó cho thấy không phải lúc nào cũng vậy, và dẫn ví dụ rằng chuyên
gia lĩnh vực có thể tự tin quá mức một cách nghiêm trọng về mức độ họ hiểu những
phương pháp XAI mà chính họ góp phần dựng nên~\autocite{p24humangap}. Công trình
được dẫn ở đó là công trình trước của chính một trong các tác giả, nên bằng
chứng duy nhất mà bài đưa ra cho việc tự tin và đúng tách nhau là bằng chứng của
chính nhóm viết bài; sách ghi lại điều đó và không rút thêm gì, vì sách không
đọc công trình được dẫn, đúng như khi
mục~\ref{sec:ch12-con-so-va-cai-no-dem} thuật ba cuộc rà trước đó qua lời bài
báo này.

\index{đánh giá!đường thứ hai trả lời câu nào|)}%
Gộp lại thì cái giá của đường thứ hai gồm hai phần và không phần nào là tiền.
Phần thứ nhất là trần đã nêu ở trên: bằng chứng mua được nói về một quan hệ
khác, nên đi hết đường này vẫn không trả lời được câu hỏi của
định nghĩa~\ref{def:ch01-do-trung-thuc}. Phần thứ hai là chất lượng: trong số ít
ỏi những lần đường này được đi, một phần đáng kể được đi theo cách mà chính bài
báo không dám gọi là bằng chứng. Hai phần ấy khác nhau về hạng. Phần thứ nhất là
giới hạn của đường, không sửa được bằng cách cố gắng hơn; phần thứ hai là chuyện
làm cẩu thả, và sửa được. Mục sau chuyển sang đường thứ ba, đường duy nhất trong
ba đường nhắm thẳng vào vế phải mà cả chương xoay quanh.
